% !TEX root = ../main.tex
\section{Qualifications for Research on Low-Cost Robot Arms in Self-Driving Labs}

\subsection{Capability}

I care about this topic because it is the work I want to pursue in graduate school. I am a
senior in mechanical engineering. In the lab I help lead a team of lab assistants building the
PiPER's enclosure, and I have run our SO-101 and PiPER arms myself. My earlier research and
product development on data acquisition and mechatronic systems taught me how measurements are
made and where their errors come from. Because the paper is a literature review, it needs three
capabilities: searching databases, which I learned in the library mini-course; reading robotics
and measurement papers closely enough to see what was actually measured, which my coursework
and data acquisition work prepared me for; and the statistics of repeatability and success
rates, which I will build during the project from \cite{chajon2026} and \cite{vincent2024}.

\subsection{Library Research Skills}\label{sec:library}

I attended the library mini-course for my major on October~27, 2026, taught by an SLC
librarian, and I will use the databases and web-searching techniques she recommended.
% CHECK: "SLC librarian" as you want it to read.
My subject resources are:
\begin{itemize}
  \item \textbf{Scopus,} a citation index of journals and conference papers in engineering,
    chemistry and materials science. Its proximity search and citation links find studies of
    low-cost arms and show who has cited them.
  \item \textbf{Engineering Village} (Compendex and Inspec), engineering indexes with controlled
    subject terms. It covers the measurement and robotics journals my AI-assisted searches
    missed, such as \emph{Measurement} and \emph{Industrial Robot}.
  \item \textbf{IEEE Xplore,} IEEE's robotics journals and conferences, including ICRA and
    \emph{IEEE Robotics and Automation Letters}, where robot benchmarks are published.
  \item \textbf{Google Scholar,} for open copies of papers and ``cited by'' links forward from my
    best sources.
  \item \textbf{Edison Scientific,} an AI agent that answers literature questions with quoted
    passages and DOIs. Two rounds of questions, on November~4 and 27, found nine of my eleven
    sources. I treat its answers as leads only, since checking them turned up misquoted passages
    and a wrong DOI.
\end{itemize}

\subsection{Review of Literature}

\textbf{Zhou \etal\ \cite{zhou2025}, ``Localization, inspection, and reasoning (LIRA) module for
autonomous workflows in self-driving laboratories.''} Zhou \etal\ study the problem at the
center of my topic: an arm that travels between stations has to find each station again when it
arrives. Their KUKA KMR iiwa mobile manipulator recalibrates at each station with a wrist camera
that reads square ArUco markers, which they compared with the older method of touching a cube.
Over 100 calibrations, a 75~mm marker located the station to within 1.02~mm and 1.01$^{\circ}$
on average, close to touch calibration's 0.92~mm and 0.90$^{\circ}$, in about 5~s instead of
53~s. In an eight-hour run between two stations, the arm completed all 600 of its pickups and
insertions.

This paper gives my list its station-level test, with an industrial baseline, and it supports
my thesis. If an arm recalibrates at every station, what matters is how well it finds a station
and repeats a short move there, not how accurately it crosses the whole lab, which is where the
arms in \cite{wattoo2024} fell short. Its weakness is that it measures with the robot's own
recorded pose rather than an independent instrument, so I will adopt its protocol and add a
dial indicator to it.

\textbf{Chaj\'{o}n \etal\ \cite{chajon2026}, ``Photogrammetric characterization of robot
positioning accuracy and repeatability.''} Chaj\'{o}n \etal\ measure an industrial
collaborative arm, a UR3, by selected elements of ISO~9283: 30 repetitions at each of five
poses, using a Raspberry Pi camera, ArUco markers on a 3D-printed tool held in the gripper, and
dial indicators as a reference. The camera measured a repeatability of 0.147 to 0.229~mm for
straight-line moves, while the dial indicators measured a mean deviation of 0.0109~mm, against 0.03~mm on the UR3's
datasheet.

This paper gives my list its arm-level test, with equipment my lab already owns, and a lesson
about measurement. The camera's numbers were more than ten times the dial indicators', so for a
UR3 the camera mostly measured its own error. A measuring tool must be finer than the arm it
measures: a camera may do for a low-cost arm with errors near a millimeter, but comparing that
arm with an industrial one takes a dial indicator, and measured and datasheet values must be
reported separately.

\subsection{Annotated Bibliography}

\begin{annotatedbib}
  \annotated{burger2020}{\RefBurger}{The landmark mobile arm that worked a lab's stations like
    a person: 688 experiments in eight days, recalibrating at every station by touching a cube,
    to a precision of $\pm$0.12~mm. It defines the role I study, and its Supplementary
    Information gives a tolerance to test against: vials and cartridges need better than
    0.5~mm.}

  \annotated{tom2024}{\RefTom}{A major review of SDLs that describes mobile robots ``operating
    multiple workstations \ldots\ mimicking a human researcher'' and prices three classes of SDL
    hardware. It frames my larger question, since it argues that open hardware can ``lower the
    financial barrier to building a SDL,'' and it supplies a counter-argument: special-purpose
    systems are more efficient than general-purpose arms.}

  \annotated{harazono2024}{\RefHarazono}{The closest study to my purpose. It tests whether a
    compact desktop arm (a Dobot M1~Pro) handles microplates accurately enough for lab
    automation, using dial gauges over 125 loaded transfers per speed, then offsets the plate on
    purpose to find how much error the task tolerates (about 2~mm). It also shows why a datasheet
    falls short: ISO~9283 ``solely considers the robot's approach to a specific location without
    holding any object.''}
  % CHECK: this quote and the 2.19 mm limit are from the bioRxiv preprint; confirm both in
  % the SLAS Technology version.

  \annotated{rupp2022}{\RefRupp}{The best evidence that a low-cost arm can do a person's lab
    work: a Dobot Magician (about 1,500 euros) pipetted, loaded an atomic absorption
    spectrometer and inoculated bacteria with ``no significant differences in accuracy and
    precision'' from a trained professional. It also feeds the counter-arguments, since the arm
    was 1.6 to 3.5 times slower and its repeatability ruled out plates with smaller wells.}

  \annotated{quigley2011}{\RefQuigley}{An early, widely cited low-cost arm whose authors report
    its performance in industrial terms (``sub-3mm repeatability,'' 1.5~m/s and a 2~kg payload)
    and explain what they gave up for price. It anchors the history of low-cost arms.}

  \annotated{doloi2025}{\RefDoloi}{A review of consumer 3D printers turned into lab equipment,
    including liquid handlers and robot arms, at 90--99\% less cost than commercial
    alternatives. It supports the larger question of democratizing SDLs and shows that a small
    lab can print its own test fixtures.}
  % CHECK: the 90-99% figure is from the article's summary; find it in the full text.

  \annotated{ferraguti2019}{\RefFerraguti}{A method for comparing arms before buying one. It
    argues that a fair comparison covers technical data with ``experimental verification of
    declared features,'' safety limits and the user's effort, and applies this to three
    collaborative arms, with collision tests under ISO/TS~15066 and measures of the effort
    needed to program them. It is my best source on the metrics that guide a purchase.}

  \annotated{wattoo2024}{\RefWattoo}{The strongest counter-argument, from inside an SDL. Putting
    two lower-precision educational arms, used in place of a UR3e, into one coordinate frame
    left an average error of 3~cm, and one arm silently ignored an infeasible command. It shows the
    failures my tests must catch.}

  \annotated{vincent2024}{\RefVincent}{It shows how to turn a few robot trials into a
    trustworthy claim with a lower confidence bound on the success rate, tighter than the common
    Clopper--Pearson bound, and which habits, such as running trials until enough succeed, make
    the bound invalid. It will set how many trials each test needs.}
\end{annotatedbib}
